% GENERATED FILE. Edit canonical lessons, then run npm run generate:books.
% canonical-source-hash: fnv1a64:3fe5ca46bade298f
% canonical-lessons: TE-C457-vasati, TE-C457-satram, TE-C457-risepsanu, TE-C457-liphtu, TE-C457-basa-ceyu

\chapter{A Place to Stay}
\label{ch:te-a-place-to-stay}
\hlchaptermodality{457}

\begin{chapteropening}
\textbf{By the end of this chapter:} \emph{I can say lodging, an inn, a reception desk, an elevator and to stay overnight.}

Complete the last lesson of chapter 457: I can say lodging, an inn, a reception desk, an elevator and to stay overnight.
\end{chapteropening}

\section[\texorpdfstring{\te{వసతి} (vasati)}{వసతి (vasati)}]{\te{వసతి} (vasati) — lodging; accommodation}

\label{lesson:TE-C457-vasati}



\begin{quote}
Before the new one: say the Telugu for dusk, then the Telugu for a stopover.
\end{quote}

\subsection*{You'll want to know: \te{వసతి}}
\textbf{\te{వసతి}} — \emph{vasati} — \textquotedblleft{}lodging; accommodation\textquotedblright{}.

\begin{quote}
\te{పండుగ} \te{రోజుల్లో} \te{హోటళ్ళలో} \te{వసతి} \te{దొరకలేదు}.
\end{quote}

\begin{quote}
\emph{paṇḍuga rōjullō hōṭaḷḷalō vasati dorakalēdu.}
\end{quote}

\begin{quote}
In the festival days we could not find lodging in the hotels.
\end{quote}

Also a hostel for students: \te{విద్యార్థి} \te{వసతి} \te{గృహం}.

\begin{cousinweb}
\textbf{\te{వసతి}} (\emph{vasati}) is borrowed from Sanskrit \emph{vasati}, \textquotedblleft{}staying, dwelling\textquotedblright{}, from \emph{vas}, \textquotedblleft{}to stay\textquotedblright{} (CDIAL 11435).
\end{cousinweb}

\begin{grammarlens}[title={What you've built}]
One more word for getting around.
\end{grammarlens}

\subsection*{Your turn}
\begin{itemize}
\raggedright
  \item \emph{Say it:} \emph{vasati}
  \item \emph{Say it:} \emph{vasati}, once more
  \item \emph{Say it:} \emph{paṇḍuga rōjullō hōṭaḷḷalō vasati dorakalēdu.} — In the festival days we could not find lodging in the hotels.
  \item \emph{Recall:} say the Telugu for dusk, then the Telugu for a stopover, then say \emph{vasati} again
\end{itemize}

\begin{tcolorbox}[breakable,colback=teal!4,colframe=teal!35!black,title={Before you move on}]
What does \te{వసతి} mean? (\textquotedblleft{}Lodging; accommodation\textquotedblright{}.) Say it once more.
\end{tcolorbox}
\section[\texorpdfstring{\te{సత్రం} (satraṁ)}{సత్రం (satraṁ)}]{\te{సత్రం} (satraṁ) — an inn; a rest house}

\label{lesson:TE-C457-satram}



\begin{quote}
Before the new one: say the Telugu for a stopover, then the Telugu for lodging.
\end{quote}

\subsection*{You'll want to know: \te{సత్రం}}
\textbf{\te{సత్రం}} — \emph{satraṁ} — \textquotedblleft{}an inn; a rest house\textquotedblright{}.

\begin{quote}
\te{పాత} \te{రోజుల్లో} \te{ప్రయాణికులు} \te{సత్రంలో} \te{ఉండేవారు}.
\end{quote}

\begin{quote}
\emph{pāta rōjullō prayāṇikulu satraṁlō uṇḍēvāru.}
\end{quote}

\begin{quote}
In the old days travellers used to stay at an inn.
\end{quote}

A free or cheap rest house for travellers and pilgrims, often run by a trust.

\begin{cousinweb}
\textbf{\te{సత్రం}} (\emph{satraṁ}) is borrowed from Sanskrit \emph{sattra}, first a long sacrificial session and later a house where food and shelter were given (Monier-Williams).
\end{cousinweb}

\begin{grammarlens}[title={What you've built}]
One more word for getting around.
\end{grammarlens}

\subsection*{Your turn}
\begin{itemize}
\raggedright
  \item \emph{Say it:} \emph{satraṁ}
  \item \emph{Say it:} \emph{satraṁ}, once more
  \item \emph{Say it:} \emph{pāta rōjullō prayāṇikulu satraṁlō uṇḍēvāru.} — In the old days travellers used to stay at an inn.
  \item \emph{Recall:} say the Telugu for a stopover, then the Telugu for lodging, then say \emph{satraṁ} again
\end{itemize}

\begin{tcolorbox}[breakable,colback=teal!4,colframe=teal!35!black,title={Before you move on}]
What does \te{సత్రం} mean? (\textquotedblleft{}An inn; a rest house\textquotedblright{}.) Say it once more.
\end{tcolorbox}
\section[\texorpdfstring{\te{రిసెప్షను} (risepṣanu)}{రిసెప్షను (risepṣanu)}]{\te{రిసెప్షను} (risepṣanu) — a reception desk}

\label{lesson:TE-C457-risepsanu}



\begin{quote}
Before the new one: say the Telugu for lodging, then the Telugu for an inn.
\end{quote}

\subsection*{You'll want to know: \te{రిసెప్షను}}
\textbf{\te{రిసెప్షను}} — \emph{risepṣanu} — \textquotedblleft{}a reception desk\textquotedblright{}.

\begin{quote}
\te{హోటలు} \te{రిసెప్షనులో} \te{తాళంచెవి} \te{ఇచ్చారు}.
\end{quote}

\begin{quote}
\emph{hōṭalu risepṣanulō tāḷañcevi iccāru.}
\end{quote}

\begin{quote}
They gave us the key at the hotel reception.
\end{quote}

The person at the desk is the \te{రిసెప్షనిస్టు}.

\begin{cousinweb}
\textbf{\te{రిసెప్షను}} (\emph{risepṣanu}) is a loan from English \emph{reception}, with the Telugu final -\te{ఉ}.
\end{cousinweb}

\begin{grammarlens}[title={What you've built}]
One more word for getting around.
\end{grammarlens}

\subsection*{Your turn}
\begin{itemize}
\raggedright
  \item \emph{Say it:} \emph{risepṣanu}
  \item \emph{Say it:} \emph{risepṣanu}, once more
  \item \emph{Say it:} \emph{hōṭalu risepṣanulō tāḷañcevi iccāru.} — They gave us the key at the hotel reception.
  \item \emph{Recall:} say the Telugu for lodging, then the Telugu for an inn, then say \emph{risepṣanu} again
\end{itemize}

\begin{tcolorbox}[breakable,colback=teal!4,colframe=teal!35!black,title={Before you move on}]
What does \te{రిసెప్షను} mean? (\textquotedblleft{}A reception desk\textquotedblright{}.) Say it once more.
\end{tcolorbox}
\section[\texorpdfstring{\te{లిఫ్టు} (liphṭu)}{లిఫ్టు (liphṭu)}]{\te{లిఫ్టు} (liphṭu) — an elevator; a lift}

\label{lesson:TE-C457-liphtu}



\begin{quote}
Before the new one: say the Telugu for an inn, then the Telugu for a reception desk.
\end{quote}

\subsection*{You'll want to know: \te{లిఫ్టు}}
\textbf{\te{లిఫ్టు}} — \emph{liphṭu} — \textquotedblleft{}an elevator; a lift\textquotedblright{}.

\begin{quote}
\te{లిఫ్టు} \te{పనిచేయలేదు}, \te{అందుకే} \te{మెట్లు} \te{ఎక్కాం}.
\end{quote}

\begin{quote}
\emph{liphṭu panicēyalēdu, andukē meṭlu ekkāṁ.}
\end{quote}

\begin{quote}
The lift was not working, so we climbed the stairs.
\end{quote}

\te{లిఫ్టు} \te{ఎక్కు}, to take the lift. In Telugu it also means a free ride: \te{లిఫ్టు} \te{ఇచ్చాడు}.

\begin{cousinweb}
\textbf{\te{లిఫ్టు}} (\emph{liphṭu}) is a loan from English \emph{lift}, with the Telugu final -\te{ఉ}.
\end{cousinweb}

\begin{grammarlens}[title={What you've built}]
One more word for getting around.
\end{grammarlens}

\subsection*{Your turn}
\begin{itemize}
\raggedright
  \item \emph{Say it:} \emph{liphṭu}
  \item \emph{Say it:} \emph{liphṭu}, once more
  \item \emph{Say it:} \emph{liphṭu panicēyalēdu, andukē meṭlu ekkāṁ.} — The lift was not working, so we climbed the stairs.
  \item \emph{Recall:} say the Telugu for an inn, then the Telugu for a reception desk, then say \emph{liphṭu} again
\end{itemize}

\begin{tcolorbox}[breakable,colback=teal!4,colframe=teal!35!black,title={Before you move on}]
What does \te{లిఫ్టు} mean? (\textquotedblleft{}An elevator; a lift\textquotedblright{}.) Say it once more.
\end{tcolorbox}
\section[\texorpdfstring{\te{బస} \te{చేయు} (basa cēyu)}{బస చేయు (basa cēyu)}]{\te{బస} \te{చేయు} (basa cēyu) — to stay overnight; to put up}

\label{lesson:TE-C457-basa-ceyu}



\begin{quote}
Before the new one: say the Telugu for a reception desk, then the Telugu for an elevator.
\end{quote}

\subsection*{You'll want to know: \te{బస} \te{చేయు}}
\textbf{\te{బస} \te{చేయు}} — \emph{basa cēyu} — \textquotedblleft{}to stay overnight; to put up\textquotedblright{}.

\begin{quote}
\te{ఆ} \te{రాత్రి} \te{మేము} \te{సత్రంలో} \te{బస} \te{చేశాం}.
\end{quote}

\begin{quote}
\emph{ā rātri mēmu satraṁlō basa cēśāṁ.}
\end{quote}

\begin{quote}
That night we put up at an inn.
\end{quote}

Past (I …): \te{బస} \te{చేశాను}, \emph{basa cēśānu}.

\te{బస}, a halt, a lodging, + \te{చేయు}. Used for a stay on a journey, away from home.

\begin{cousinweb}
\textbf{\te{బస} \te{చేయు}} (\emph{basa cēyu}) is a Telugu phrase: \te{బస}, \textquotedblleft{}a halt, a lodging\textquotedblright{}, whose origin is not settled, and the native \te{చేయు}, \textquotedblleft{}to do\textquotedblright{} (DEDR 1957).
\end{cousinweb}

\begin{grammarlens}[title={What you've built}]
One more word for getting around.
\end{grammarlens}

\subsection*{Your turn}
\begin{itemize}
\raggedright
  \item \emph{Say it:} \emph{basa cēyu}
  \item \emph{Say it:} \emph{basa cēyu}, once more
  \item \emph{Say it:} \emph{ā rātri mēmu satraṁlō basa cēśāṁ.} — That night we put up at an inn.
  \item \emph{Recall:} say the Telugu for a reception desk, then the Telugu for an elevator, then say \emph{basa cēyu} again
\end{itemize}

\begin{tcolorbox}[breakable,colback=teal!4,colframe=teal!35!black,title={Before you move on}]
What does \te{బస} \te{చేయు} mean? (\textquotedblleft{}To stay overnight; to put up\textquotedblright{}.) Say it once more.
\end{tcolorbox}
